\section{The explicit robustness constant}\label{sec:explicit-robust}

The lower comparison of Ref.~\cite{TG} uses a flagged-sector estimate from
Ref.~\cite{PP}. We obtain explicit constants for the sector estimate and
the nine-copy comparison. Write $\operatorname{Sch}(v)$ for the decreasing
Schmidt-amplitude list of an arbitrary vector $v$, padded by zeroes. Write
$k\times s$ for $k$ copies of a list $s$, arranged in decreasing order.
We use
\begin{equation}\label{eq:sch-contract-explicit}
 \|\operatorname{Sch}(v)-\operatorname{Sch}(w)\|_2\le\|v-w\|.
\end{equation}
This follows by applying von Neumann's trace inequality to the
Hilbert--Schmidt coefficient operators of the two vectors and expanding
both squares. Finite-rank approximation gives the same inequality in
arbitrary Hilbert spaces.

\begin{lemma}[Quantitative flagged-sector estimate]\label{lem:explicit-sector}
Let $\psi$ be a vector and the measurements of one flagged magic-square
copy be projective. Let $F$ be Alice's $R_1$ flag projection and let $N$
be the sum of the forbidden masses on the 25 internal pairs, evaluated
on $\psi$. Then
\begin{equation}\label{eq:explicit-sector}
 \|\operatorname{Sch}((I-F)\psi)
      -4\times(\sqrt2\,\operatorname{Sch}(F\psi))\|_2\le 210\sqrt N .
\end{equation}
\end{lemma}

\begin{proof}
We omit tensor identities for local operators. Denote Alice's row flag
projections by $F_i$, and Bob's column flag projections by $G_j$. Put
$F=F_1$, $G_\star=G_1$, and $\mu=(I-F)(I-G_\star)\psi$. For the row--column
pair $(i,j)$ let $u_{ij}^2$ be its forbidden mass and let
$M=\sum_{ij}u_{ij}^2$. Define
\[
 d_A=\|F(I-G_\star)\psi\|,\qquad
 d_B=\|(I-F)G_\star\psi\|,\qquad d^2=d_A^2+d_B^2\le u_{11}^2.
\]
The last inequality and $\|(F_i-G_j)\psi\|\le u_{ij}$ use the flag-mismatch
rule.

\paragraph{Extracting the magic-square state.}
For each cell let $A_{ij},B_{ij}$ be the two local sign observables, giving
every flag the sign $+1$. They are self-adjoint unitaries and satisfy
$\|(A_{ij}-B_{ij})\psi\|\le\beta_{ij}$, where $\beta_{ij}=2u_{ij}$.
Alice's rows and Bob's columns commute within each context. Alice's row
products and Bob's first two column products are $I$, and
\[
 B_{33}=-B_{13}B_{23}+2G_3.
\]
Denote the generators at cells $(11,12,21,22)$ by $(a,b,c,d_0)$ on Alice
and by the same primed symbols on Bob. A word of $k$ local observables on
$\psi$ differs from the reversed word on the other party by at most the
sum of the $k$ cell bounds $\beta_{ij}$. To see this, replace the
rightmost factor on $\psi$ and commute it to the left of the remaining
factors. In particular, a difference $L$ of two signed local words,
followed by a one-letter suffix $s$, has
$\|Ls\psi\|\le\|L\psi\|+2\beta_s$.

Write $W=\sum_{ij}\beta_{ij}$, and let $[q,r]=qr-rq$ and
$\{q,r\}=qr+rq$. On either party the three expressions
\[
 d_0\{b,c\}a,\qquad b\{a,d_0\}c,\qquad c\{a,d_0\}b
\]
applied to $\psi$ differ from $2G_3\psi$ by at most $W$.
In each expression the two terms are a product of column pairs and a
product of row pairs. The column pairs are the cells $31,32$ on Bob,
and their product transfers to Alice's $A_{33}$. The row pairs are the
cells $13,23$ on Alice, and their product transfers to Bob's
$B_{13}B_{23}$. The remaining transfer is the cell $33$.
Expanding each product on the other party by the reversed-word rule
costs in total each of the nine bounds $\beta_{ij}$ once. Reversing the
order of the two row or column pairs is allowed at the transfer, by
commutation in Bob's column 3 or Alice's row 3. Also, if two generators
$q,r$ commute on the other party, then
$\|[q,r]\psi\|\le2(\beta_q+\beta_r)$ by two word transfers.

We use the real two-qubit Pauli group $\mathcal P$ generated by
$a=ZI$, $b=IZ$, $c=IX$, $d_0=XI$ and $J=-I$, with the usual Pauli
matrices $Z,X$. Its 32 elements have unique signed normal words in any
fixed ordering of $a,b,c,d_0$, with each generator appearing at most once.
The anticommuting pairs are $(a,d_0)$ and $(b,c)$, and the other pairs
commute. For $g\in\mathcal P$, let $A_g$ evaluate its signed normal word
in the physical observables, using the order $(c,d_0,a,b)$. Define
$B_g$ on Bob using $(b,d_0,a,c)$. The local isometries
\[
 V_A v=32^{-1/2}\sum_{g\in\mathcal P}|g\rangle A_g v,\qquad
 V_B v=32^{-1/2}\sum_{g\in\mathcal P}|g\rangle B_g v
\]
have their registers in the $-1$ eigenspace of $\rho(J)$, where
$\rho(h)|g\rangle=|gh^{-1}\rangle$ is the right regular action.
The exact commutation of the last two generators in each order gives
$V_A a=\rho(a)V_A$, $V_A b=\rho(b)V_A$, and
$V_B a'=\rho(a)V_B$, $V_B c'=\rho(c)V_B$ on the full input spaces.

For a generator $h$ write $\alpha_{Ah}=\|(V_A A_h-\rho(h)V_A)\mu\|$,
and define $\alpha_{Bh}$ in the same way. These norms are the
root-mean-square errors of rewriting the physical word for $g$
followed by $h$ into the word for $gh$. Put
\[
 L_A=W+2u_{13}+2d_A,\qquad L_B=W+2u_{13}.
\]
The table gives the rewriting bounds on $(I-G_\star)\psi$ for Alice and
$(I-F)\psi$ for Bob. The three columns cover presence of only the
first, only the second, and both of the last two letters in the word
for $g$. With neither letter present the rewriting is exact.
\[
\begin{array}{c|ccc}
 &\text{first}&\text{second}&\text{both}\\ \hline
 A,c
 &2(\beta_a+\beta_c)& L_A+2\beta_a
 & L_A+2\beta_a+2(\beta_a+\beta_c+\beta_b)\\
 A,d_0
 &L_A+2\beta_c&2(\beta_b+\beta_{d_0})
 &L_A+2(\beta_b+\beta_{d_0})\\
 B,b
 &2(\beta_a+\beta_b)&L_B+2\beta_a
 &L_B+2\beta_a+2(\beta_a+\beta_b+\beta_c)\\
 B,d_0
 &L_B+2\beta_b&2(\beta_c+\beta_{d_0})
 &L_B+2(\beta_c+\beta_{d_0})
\end{array}
\]
For the anticommutator swaps the flag correction costs at most
$2\|(I-G_\star)G_3\psi\|\le2u_{13}+2d_A$ on Alice and
$2\|(I-F)G_3\psi\|\le2u_{13}$ on Bob. When the swap has no suffix,
transfer the trailing letter of the displayed anticommutator identity
to the other party before projecting. Use $a$ for the pair $(b,c)$, $c$
on Alice for $(a,d_0)$, and $b$ on Bob for $(a,d_0)$. The transferred
letter commutes with the opposite flag $G_\star$ or $F$. When the
$(a,d_0)$ swap already has a suffix, use the identity with that suffix
and project directly. In the remaining swaps only the commutator
bound and, if present, one suffix transfer are needed.

On $\mu$ add $2d_A$ or $2d_B$ to each nonzero case in the table, because
the difference of the two signed words has norm at most two. Each case
has register weight $1/4$. Write
$v=(u_{11},u_{12},\ldots,u_{33},d_A,d_B)^{\mathsf T}$ in row-major
order, and let $L_{Ah}$ and $L_{Bh}$ be the coefficient matrices of
the three bounds for each row of the table, including the additions
$2d_A$ or $2d_B$. Thus $\alpha_{Ah}\le\|L_{Ah}v\|_2/2$ and
likewise for Bob.

The consistency errors $\delta_h=\|(A_h-B_h)\mu\|$ satisfy
\[
 \delta_a\le 2u_{11},\quad
 \delta_b\le 2u_{12}+2d_B,\quad
 \delta_c\le 2u_{21}+2d_A,\quad
 \delta_{d_0}\le 2u_{22}+u_{12}+u_{21}+d_A+d_B .
\]
For $a,b,c$, use commutation with the corresponding anchor flags.
For $d_0$, the cell difference commutes with $G_2$ and with $F_2$.
Compare $\mu$ to $(I-G_2)\psi$ or $(I-F_2)\psi$, at distances at most
$d_B+u_{12}$ and $d_A+u_{21}$, respectively, and average the two bounds.
Write $\ell_h(v)$ for the corresponding linear upper bound on $\delta_h$.

The $-1$ register sector carries $ZI,IZ,IX,XI$ on two qubits and an
identity on a multiplicity space. Indeed the four joint eigenspaces of
$\rho(a),\rho(b)$ are identified with the joint $+1$ eigenspace by
$\rho(d_0),\rho(c)$, which flip their respective signs and commute.
Let $w=(V_A\otimes V_B)\mu$. The four simultaneous involutions
$K_h=\rho_A(h)\rho_B(h)$ commute and their joint $+1$ projection
$\Pi$ has range $\Phi_4\otimes\mathcal K_A\otimes\mathcal K_B$.
The generators $ZI,IZ$ first enforce equal bits on both registers and
$IX,XI$ then enforce equal coefficients on the four equal-bit basis
vectors. Put $w_0=\Pi w=\Phi_4\otimes j$ and $\eta=\|w-w_0\|$. Since
$\|(I-K_h)w\|\le\alpha_{Ah}+\delta_h+\alpha_{Bh}$,
\[
 4\eta^2\le\sum_h(\alpha_{Ah}+\delta_h+\alpha_{Bh})^2\le v^{\mathsf T}Hv .
\]
The first inequality follows from $I-\Pi\le\sum_h(I-K_h)/2$.
For the second, use weighted Cauchy--Schwarz with weights $(1,6)$
for $b$, $(1,5)$ for $c$, and $(2,6,5)$ for $d_0$. In terms of the
linear forms above the resulting matrix is given by
\begin{align*}
 v^{\mathsf T}Hv
 &=\ell_a(v)^2+7\ell_b(v)^2+6\ell_c(v)^2+\tfrac{13}{2}\ell_{d_0}(v)^2\\
 &\quad+\tfrac7{24}\|L_{Bb}v\|_2^2+\tfrac3{10}\|L_{Ac}v\|_2^2
   +\tfrac{13}{24}\|L_{Ad_0}v\|_2^2+\tfrac{13}{20}\|L_{Bd_0}v\|_2^2.
\end{align*}
All its entries are nonnegative. For
\[
 D_*=\operatorname{diag}(.561,1,1,1,1,1,1,1,1,.439,.439),\qquad
 q=(1000,623,452,634,517,226,226,226,226,833,553)^{\mathsf T},
\]
the integers $\lceil1000(Hq)_i/(D_*q)_i\rceil$ in order are
\[
 (521543,521579,521308,521552,521430,521308,521308,521308,521308,521546,521180).
\]
Thus $Hq\le521.6D_*q$. Using
$2v_iv_j\le v_i^2q_j/q_i+v_j^2q_i/q_j$ entrywise gives
$v^{\mathsf T}Hv\le521.6v^{\mathsf T}D_*v\le521.6M$, the last
inequality by $d_A^2+d_B^2\le u_{11}^2$. This gives
$\eta\le11.42\sqrt M$.

\paragraph{Comparing the flag and magic-square sectors.}
Embed each local space into a direct sum of a flag space and an
extracted magic-square space, by $v\mapsto Fv\oplus V_A(I-F)v$ for Alice
and $v\mapsto G_\star v\oplus V_B(I-G_\star)v$ for Bob. Let $P,Q$ be the two
flag-space projections and let $\Psi$ be the embedded state. Write
$\varphi=PQ\Psi$, $\zeta=\varphi+w_0$, and $e=\Psi-\zeta$. The components
of $e$ in $P(I-Q)$, $(I-P)Q$ and $(I-P)(I-Q)$ have norms
$d_A,d_B,\eta$, respectively. The measurements transport to the images
of the isometries and can be extended as projective measurements by
assigning the orthogonal complement to one outcome.

On either magic register let $R_y$ be the projection onto
$r_y=|0\rangle|y\rangle$ and $C_x$ the projection onto
$c_x=|0\rangle(|0\rangle+(-1)^x|1\rangle)/\sqrt2$, tensored with the
identity on its multiplicity and original local space. Extend these
projections by zero on the flag spaces and use superscripts $A,B$
for the party. The exact intertwining for Alice's $a,b$ and Bob's
$a',c'$ means that $P_y=P+R_y^A$ and $Q_x=Q+C_x^B$ act on $\Psi$
as the calibrated physical coarse-grainings $f\ \text{or}\ r_y$
and $f\ \text{or}\ c_x$. Put $E_x^0=P+C_x^A$ and $H_y^0=Q+R_y^B$.
These satisfy $E_x^0\zeta=Q_x\zeta$ and $H_y^0\zeta=P_y\zeta$.

Let $E_x,H_y$ be the actual sign projections of $T_x,U_y$, and let
$t_x,u_y'$ be the forbidden masses at $(T_x,C_1)$ and $(R_1,U_y)$.
The mismatch rule gives
$\|(E_x-Q_x)\Psi\|\le\sqrt{t_x}$ and
$\|(H_y-P_y)\Psi\|\le\sqrt{u_y'}$. Set $T=\sum_x t_x+\sum_y u_y'$.
Define $\xi_{xy}=E_x^0H_y^0\zeta$. Using the two exact relations on
$\zeta$ gives
\begin{align*}
 E_xH_y\Psi-\xi_{xy}
 &=H_y(E_x-Q_x)\Psi+E_x^0(H_y-P_y)\Psi\\
 &\quad+H_y(Q_x-E_x^0)e+E_x^0P_y e .
\end{align*}
The first two terms together have norm at most $2\sqrt T$ in the
direct sum over $x,y$. For the third, each difference of projections
has operator norm at most one, so the direct-sum norm is at most
$2\sqrt{d^2+\eta^2}$. For the fourth, $E_x^0P_y$ is the identity on
Alice's flag and equals $C_x^A R_y^A$ on her magic space, with
$\sum_{xy}R_y^A C_x^A R_y^A\le I$. Hence
\begin{equation}\label{eq:explicit-sign-support}
 \left(\sum_{xy}\|E_xH_y\Psi-\xi_{xy}\|^2\right)^{1/2}
 \le 2\sqrt T+2\sqrt{d^2+\eta^2}+\sqrt{\eta^2+4d_A^2+d_B^2}.
\end{equation}

Let $U_x,V_y$ be the local sign unitaries, acting by the measured sign
on the sign subspace and by $+1$ on its complement. Put
$t_{xy}=(-1)^{xy}$ and let $S$ sum the forbidden masses on the four
sign pairs. Then
\[
 \sum_{xy}\|(U_x-t_{xy}V_y)E_xH_y\Psi\|^2=4S.
\]
The ideal support vector is $\xi_{xy}=\varphi+\chi_{xy}$ with
$\chi_{xy}=t_{xy}|c_xr_y\rangle h$ and $h=j/(2\sqrt2)$, by the
overlap of the two selected magic-register vectors. Define
$b_{xy}=(U_x-t_{xy}V_y)\xi_{xy}$ and $B^2=\sum_{xy}\|b_{xy}\|^2$.
Equation~\eqref{eq:explicit-sign-support} gives
\[
 B\le 2\sqrt S+4\sqrt T+4\sqrt{d^2+\eta^2}
                 +2\sqrt{\eta^2+4d_A^2+d_B^2}.
\]

To compare $\varphi$ and $h$, project the four $b_{xy}$ onto $PQ$.
The resulting differences are $P U_x\varphi-t_{xy}Q V_y\varphi$.
For four vectors $v_x,w_y$ in a common space,
$2w_0=(w_0-v_0)+(v_0-w_1)+(w_1+v_1)+(w_0-v_1)$.
Cauchy--Schwarz therefore bounds $\|w_0\|^2$ by the sum of the four
squared differences $\|v_x-t_{xy}w_y\|^2$, and the same holds for
$v_0$. In particular,
$\|Q V_0\varphi\|^2\le\sum_{xy}\|PQ b_{xy}\|^2$.
Next project onto $C_x^A R_y^B$ and remove the explicit qubit factors
and the scalar sign $t_{xy}$. The four differences now involve the
compressions of $U_x$ to $c_x$ and of $V_y$ to $r_y$, acting on the
common vector $h$. This gives
$\|C_0^A U_0\chi_{00}\|^2\le\sum_{xy}\|C_x^A R_y^B b_{xy}\|^2$.

The components of $U_0\chi_{00}-V_0\varphi$ in $PQ$ and $C_0^A$
have exactly these two bounded norms. Its remaining components, in
$P(I-Q)$ and $I-P-C_0^A$, are
$P(I-Q)b_{00}$ and $(I-P-C_0^A)R_0^B b_{00}$, respectively.
These are orthogonal to one another and to the two projections
of $b_{00}$ used in the cycle bounds. Adding the squared norms yields
$\|U_0\chi_{00}-V_0\varphi\|\le B$.
By~\eqref{eq:sch-contract-explicit} and local unitarity,
$\|\operatorname{Sch}(h)-\operatorname{Sch}(\varphi)\|_2\le B$.

The list of $w_0$ is $4\times(\sqrt2\,\operatorname{Sch}(h))$.
The extracted $F\psi$ is within $d_A$ of $\varphi$, and the extracted
$(I-F)\psi$ is within $\sqrt{d_B^2+\eta^2}$ of $w_0$. The left side of
~\eqref{eq:explicit-sector} is therefore bounded by
\[
 \sqrt{d_B^2+\eta^2}+2\sqrt2\,d_A+2\sqrt2\,B
 \le 209.6\sqrt M+8\sqrt2\sqrt T+4\sqrt2\sqrt S .
\]
For the numerical bound use
$\sqrt{1+11.42^2}<11.464$, $\sqrt{4+11.42^2}<11.594$ and
$2\sqrt2<2.828428$. Since $209.6^2+160<210^2$ and $M+T+S=N$,
Cauchy--Schwarz proves~\eqref{eq:explicit-sector}.
\end{proof}

\begin{proof}[Lower bound in Theorem~\ref{thm:equiv}]
First make the measurements projective by local isometries. For a
finite POVM $(M_a)_a$ on $\mathcal H$, its isometry
$Wv=\sum_a\sqrt{M_a}v\otimes|a\rangle$ maps into an outcome register
space $\mathcal H\otimes\mathbb C^m$. A common enlarged space is
$\mathcal H\oplus(\mathcal H\otimes\mathbb C^m)$, with the input in the
first summand and $m$ at least the largest number of outcomes. For
each question the self-adjoint unitary
$(v,w)\mapsto(W^*w,Wv+(I-WW^*)w)$ exchanges the input with the range
of $W$. Conjugating the outcome-coordinate projections, with the
first summand and any unused coordinates assigned to a fixed
outcome, gives the required projective measurement. This dilation
preserves correlations and the nonzero Schmidt amplitudes.

Write $\Lambda=\operatorname{Sch}(\psi)$, $Q_j$ for Bob's label
projections, and $F_j$ for the nine $R_1$ flags on Alice. Let
$d_j=\|(F_j-Q_j)\psi\|$, $N_L=\sum_j d_j^2$, and $N_F=\sum_j N_j$,
where $N_j$ is the internal forbidden mass of copy $j$. For the
strategy's losing probability $\varepsilon$ the disjoint budgets
give $N_F+N_L\le243\varepsilon$.

To separate the label sectors, use the local column operator
$Tv=\sum_j F_jv\otimes|j\rangle$ on Alice, with
$A=T^*T=\sum_jF_j$. Add an unused register entry and extend the partial
isometry in the polar decomposition $T=V A^{1/2}$ to an isometry
$V$ by mapping $\ker A$ into that entry. Since $A\ge F_j$, operator
monotonicity of the square root gives $A^{1/2}\ge F_j$. On the
orthogonal Bob labels this yields $A^{1/2}\ge\sum_j F_jQ_j$, and
\[
 \|(V-T)\psi\|^2=\|(I-A^{1/2})\psi\|^2
 \le \langle\psi,(A+I-2\textstyle\sum_j F_jQ_j)\psi\rangle=N_L.
\]
Apply the exact label isometry on Bob. For any label subset $I_0$,
the raw map $T$ on $\sum_{j\in I_0}Q_j\psi$ then differs from
$\sum_{j\in I_0}Q_j\psi\otimes|j,j\rangle$ by a vector of squared
norm $\sum_{j\in I_0,k}\|Q_j(F_k-Q_k)\psi\|^2\le N_L$.
The same input restricted to $I_0$ can only decrease the
$(V-T)$ error. Thus the two local isometries give a vector within
$2\sqrt{N_L}$ of the target.
Let $S_j=\operatorname{Sch}(Q_j\psi)$ and let $\bigsqcup$ denote
sorted union. The target is orthogonal on both sides, so
\[
 \left\|\operatorname{Sch}\left(\sum_{j\in I_0}Q_j\psi\right)
      -\bigsqcup_{j\in I_0}S_j\right\|_2\le2\sqrt{N_L}.
\]

Use this estimate on each complementary label set and replace $F_j$
by $Q_j$ in~\eqref{eq:explicit-sector} using
~\eqref{eq:sch-contract-explicit}. Then
\[
 \left\|\bigsqcup_{k\ne j}S_k-4\times(\sqrt2 S_j)\right\|_2
 \le 210\sqrt{N_j}+(1+2\sqrt2)d_j+2\sqrt{N_L}.
\]
Write $U=\bigsqcup_j S_j$ and
$\Delta(s)=\|2\times s-\sqrt2\,s\|_2$. Concatenate the nine list
comparisons and sort each side. Sorting cannot increase the
$\ell^2$ distance, by the rearrangement inequality. Each $S_k$
appears eight times on the left and four times on the right, giving
\[
 2\Delta(U)\le210\sqrt{N_F}+(7+2\sqrt2)\sqrt{N_L}.
\]
The all-label estimate gives $\|\Lambda-U\|_2\le2\sqrt{N_L}$.
Since $\Delta$ is $2\sqrt2$-Lipschitz,
\[
 \Delta(\Lambda)\le105\sqrt{N_F}+c_L\sqrt{N_L},\qquad
 c_L=\frac72+5\sqrt2,\qquad c_L^2<112 .
\]
By Lemma~\ref{lem:formula}, $\cE(\psi)=\Delta(\Lambda)^2/2$.
Cauchy--Schwarz and the disjoint budgets now give
\[
 \cE(\psi)\le \frac{243(210^2+448)}{8}\,\varepsilon
              =1{,}353{,}145.5\,\varepsilon .
\]
This holds for every strategy and hence also for
$\varepsilon=1-\omega(\psi)$.
\end{proof}
